\section*{6 Bit pigeonhole principle}
\subsection*{6.1 Reduction from $\mathrm{BPHP}\circ\oplus_k$ to $\mathrm{BPHP}$}
Let $T\subseteq X_1\times\cdots\times X_k\times Y$ and $S\subseteq Z_1\times\cdots\times Z_k\times W$ be relations. We say $S$ is many-one reducible to $T$ if there are maps $f_i:X_i\to Z_i$ and $g:W\to Y$ such that if $(f_1(x_1),\ldots,f_k(x_k),y)\in T$ then $(x_1,\ldots,x_k,g(y))\in S$.
\begin{authorclaim}{Lemma39}
If $S$ is many-one reducible to $T$, then $R^{\mathrm{pub}}_{1/3}(S)\le R^{\mathrm{pub}}_{1/3}(T)$.
\end{authorclaim}
\begin{proof}
The $i$th party computes $f(x_j)$ for all $j\in[k]\setminus\{i\}$ and then all parties run the optimal protocol for $T$. As soon as all the parties learn an answer $y$ they compute $g(y)$ without communication.
\end{proof}
Recall $\mathrm{BPHP}_{2^n}^M$ encodes $M$ different strings $s_1,\ldots,s_M\in\{0,1\}^n$. For $n=\ell k+r$, $0\le r<k$, partition each row $s_i=s_i^{(1)}\cdots s_i^{(k)}$, where $|s_i^{(t)}|=\ell+1$ for $t\le r$ and $\ell$ for $t>r$. Partition $\Pi_k$ puts all $s_1^{(t)},\ldots,s_M^{(t)}$ in part $t$.
$\operatorname{SearchPair}_{2^n}^M$ asks for distinct $i,j$ with $s_i=s_j$ from variables partitioned according to $\Pi_k$.
\begin{authorclaim}{Proposition40}
$\operatorname{SearchPair}_{2^n}^M$ is many-one reducible to $\operatorname{Search}(\mathrm{BPHP}_{2^n}^M)$ with partition $\Pi_k$.
\end{authorclaim}
\begin{proof}
The proof is straightforward.
\end{proof}
\begin{authorclaim}{Theorem41}
$\oplus_k\mathrm{BPHP}_{2^\ell}^m$ is many-one reducible to $\operatorname{SearchPair}_{2^{k\ell}}^{m2^{(k-1)\ell}}$.
\end{authorclaim}
\begin{proof}
Let us denote $M=m\cdot2^{(k-1)\ell}$. Consider a set
$Z=\{(y_1,y_2,\ldots,y_k)\in(\mathbb F_2^\ell)^k\mid\sum_iy_i=0\}$.
It is easy to see that $|Z|=2^{(k-1)\ell}$. Let $\phi$ be a bijection between $[M]$ and $Z\times[m]$.
Let for $i\in[m]$ and $t\in[k]$, $x_i^{(t)}$ denote the $i$th string of the $t$th party in the communication problem $\oplus_k\mathrm{BPHP}_{2^\ell}^m$. Let $x_i:=(x_i^{(1)},\ldots,x_i^{(k)})$.
For every $t\in[k]$ we define $f_t$ as follows: $f_t(x_1^{(t)},\ldots,x_m^{(t)})$ is a sequence of rows $r_1^{(t)},r_2^{(t)},\ldots,r_M^{(t)}$ such that for all $i\in[M]$, $r_i^{(t)}=z_t+x_j^{(t)}$, where $(z,j)=\phi(i)$ for all $z\in Z$ and $j\in[m]$ (recall that $z\in Z$ is divided on $k$ parts of equal lengths and $z_t$ denotes the $t$th part).
Let us construct the function $g$ from the definition of the reduction.
Let $q,w\in[M]$ and $q\ne w$. Assume that $\phi(q)=(z,j_1)$ and $\phi(w)=(z,j_2)$. We define $g(q,w):=(j_1,j_2)$.
Let us verify that $f_1,f_2,\ldots,f_k$ and $g$ define a reduction. Let $q,w\in M$ be a pair of different numbers such that the assignment $\alpha:=\{s_i\leftarrow r_i^{(1)}r_i^{(2)}\cdots r_i^{(k)}\mid i\in[M]\}$ satisfies $s_q=s_w$.
Assume that $g(q,w)=(j_1,j_2)$. We need to verify that $j_1\ne j_2$ and $\sum_{t=1}^kx_{j_1}^{(t)}=\sum_{t=1}^kx_{j_2}^{(t)}$.
Notice that under the assignment $\alpha$ the value of $s_q$ is $x_{j_1}+z$ and the value of $s_w$ is $x_{j_2}+y$, where $j_1,j_2\in[m]$ and $z,y\in Z$ such that $(z,j_1)=\phi(q)$ and $(y,j_2)=\phi(w)$. If $j_1=j_2$, then $x_{j_1}+z=x_{j_2}+y$ implies $z=y$. Since $\phi$ is a bijection, we get $q=w$. Thus, $j_1\ne j_2$.
For each $t\in[k]$, the following equality holds.
\begin{equation}
z_t+x_{j_1}^{(t)}=y_t+x_{j_2}^{(t)}.\tag{3}
\end{equation}
If we sum up equations(3) for all $t\in[k]$ and use that $y,z\in Z$, we get $\sum_{t=1}^kx_{j_1}^{(t)}=\sum_{t=1}^kx_{j_2}^{(t)}$. Hence, $(j_1,j_2)$ is a correct answer for $\oplus_k\mathrm{BPHP}_{2^\ell}^m$.
\end{proof}
The following proposition deals with the case, where the number of bits is not divisible by $k$.
\begin{authorclaim}{Proposition42}
For $n=k\ell+r$, $0\le r<k$, $M>2^{k\ell}$, $\operatorname{SearchPair}_{2^{k\ell}}^M$ is many-one reducible to $\operatorname{SearchPair}_{2^n}^{M2^r}$.
\end{authorclaim}
\begin{proof}
Let $x_1,x_2,\ldots,x_M$ be the input of $\operatorname{SearchPair}_{2^{k\ell}}^M$, let $x_j^{(t)}$ be the $t$th part of the row $x_j$ according to the partition $\Pi_k$. Given this input we construct an input for $\operatorname{SearchPair}_{2^n}^{M2^r}$.
Let $\tau$ be a bijection between $[M]\times\{0,1\}^r$ and $[M2^r]$.
For each $i\in[M]$ we construct $2^r$ rows $y_{\tau(i,\alpha)}$ one for each $\alpha\in\{0,1\}^r$. Let $\Pi_k$ partition a row $y_{\tau(i,\alpha)}$ into the following parts: $y_{\tau(i,\alpha)}^{(1)}y_{\tau(i,\alpha)}^{(2)}\cdots y_{\tau(i,\alpha)}^{(k)}$. Let
\[
y_{\tau(i,\alpha)}^{(t)}
=\begin{cases}x_i^{(t)},&t>r,\\x_i^{(t)}\alpha_t,&0\le t\le r.\end{cases}
\]
Now we can define the function $f_t(x_1^{(t)},\ldots,x_M^{(t)})$ as $y_{\tau(i,\alpha)}^{(t)}$ for each $i\in M$ and $\alpha\in\{0,1\}^r$ and $t\in[k]$. Observe that for each $i\in[M]$ the rows $y_{i,\alpha}$ for $\alpha\in\{0,1\}^r$ are distinct. That allows us to define the function $g$ as $g(\tau(i_1,\alpha_1),\tau(i_2,\alpha_2))=(i_1,i_2)$. All the required properties can be easily verified.
\end{proof}
\begin{authorclaim}{Theorem4}
Let $M=2^n+2^{k+n-\lfloor n/k\rfloor}$ and $n\ge k(k+7)$. With partition $\Pi_k$,
\[
R^{\mathrm{pub}}_{1/3}(\operatorname{Search}(\mathrm{BPHP}_{2^n}^M))
=\Omega(2^{n/(2k)-3k/2}/k).
\]
For $k=2$ the stronger bound is $\Omega(2^{n/2})$.
\end{authorclaim}
\begin{proof}
Let $\ell=\lfloor n/k\rfloor$, $r=n-\ell k$. By Proposition40, Proposition42, Theorem41 and Corollary19, respectively,
\begin{align*}
R^{\mathrm{pub}}_{1/3}(\operatorname{Search}(\mathrm{BPHP}_{2^n}^M))
&=R^{\mathrm{pub}}_{1/3}(\operatorname{Search}(\mathrm{BPHP}_{2^n}^{(2^k+2^\ell)2^{(k-1)\ell+r}}))\\
&\ge R^{\mathrm{pub}}_{1/3}(\operatorname{SearchPair}_{2^n}^{(2^k+2^\ell)2^{(k-1)\ell+r}})\\
&\ge R^{\mathrm{pub}}_{1/3}(\operatorname{SearchPair}_{2^{k\ell}}^{(2^k+2^\ell)2^{(k-1)\ell}})\\
&\ge R^{\mathrm{pub}}_{1/3}(\oplus_k\mathrm{BPHP}_{2^\ell}^{2^\ell+2^k})\\
&=\Omega(2^{\ell/2-3k/2}/k)=\Omega(2^{n/(2k)-3k/2}/k).
\end{align*}
The case of $k=2$ can be treated in the same way, the only difference is in the application of Corollary19.
\end{proof}
\subsection*{6.2 Upper bound for communication complexity}
\begin{authorclaim}{Proposition5}
For $M>2^n$ and $k\in\{2,\ldots,n\}$ there is a deterministic NOF protocol for $\operatorname{Search}(\mathrm{BPHP}_{2^n}^M)$ with partition $\Pi_k$ transmitting $O(2^{\lceil n/k\rceil}\log M)$ bits.
\end{authorclaim}
\begin{proof}
The protocol is going to have only two active parties: the second party, which we call Alice, and the first party, which we call Bob. We are going to use that Alice can see the variables $s_1^{(1)},\ldots,s_M^{(1)}$ and that Bob can see all other variables.
Let us denote $\bar s_i=s_i^{(1)}s_i^{(2)}s_i^{(3)}\cdots s_i^{(k)}\in\{0,1\}^{n-\lceil n/k\rceil}$ the bits Bob sees in the $i$th line for $i\in[M]$. Bob finds a value $\alpha\in\{0,1\}^{n-\lceil n/k\rceil}$ such that the size of the set $S_\alpha=\{i\in[M]\mid\bar s_i^{(1)}=\alpha\}$ is larger than $2^{\lceil n/k\rceil}$. Such $\alpha$ exists since $M>2^n$. Bob then picks an arbitrary subset $S'$ of $S_\alpha$ of size $2^{\lceil n/k\rceil}+1$ and sends the description of $S'$ to Alice using $(2^{\lceil n/k\rceil}+1)\cdot\lceil\log_2M\rceil$ bits. Then, by the pigeonhole principle there exists $i\ne j\in S'$ such that $s_i^{(1)}=s_j^{(1)}$. Alice and Bob then spend $O(\log M+n)$ bits transmitting indices $i$ and $j$ and all the values of the $i$th and $j$th lines to each other. Both of them then find the falsified clause of $\mathrm{BPHP}_{2^n}^M$ with no communication because it only depends on variables $s_i$ and $s_j$ and broadcast its description to all of the parties using an additional $O(n+\log M)$ bits.
\end{proof}
For $k=2$ this coincides with the lower bound up to a logarithmic factor. For larger $k$ they are polynomially related; the dependence on $k$ is a genuine phenomenon.
\subsection*{6.3 Short $\mathrm{Th}(\log n)$ proof of $\mathrm{BPHP}_n^m$}
This gives a short tree-like $\mathrm{Th}(\log n)$ refutation, similar to the conversion of unary PHP resolution proofs to $\mathrm{Res}(\log n)$ proofs of binary PHP in [5].
\begin{authorclaim}{Proposition43}
If there exists a tree-like $\mathrm{Th}(1)$ refutation of $\mathrm{PHP}_{2^\ell}^m$ of size $S$, there exists a tree-like $\mathrm{Th}(\ell)$ refutation of $\mathrm{BPHP}_{2^\ell}^m$ of size $O(S)$.
\end{authorclaim}
\begin{proof}
Let $p_{i,j}$ for $i\in[m]$ and $j\in[2^\ell]$ be a variable of $\mathrm{PHP}_{2^\ell}^m$ indicating that the $i$th pigeon flies to the $j$th hole. Let $s_{i,k}$ for $i\in[m]$, $k\in[\ell]$ be a variable of $\mathrm{BPHP}_{2^\ell}^m$ indicating the $\ell$th bit of the $i$th string $s_i$.
Let $Q_j(x_1,x_2,\ldots,x_\ell)$ for $j\in[2^k]$ be a multilinear polynomial over reals such that for all $a_1,a_2,\ldots,a_\ell\in\{0,1\}^\ell$, $Q_j(a_1,a_2,\ldots,a_\ell)=1$ if $(a_1,a_2,\ldots,a_\ell)=\operatorname{bin}_\ell(j-1)$ and $Q_j(a_1,a_2,\ldots,a_\ell)=0$ otherwise. We ma define $Q_j$ as follows
\[
Q_j(x_1,\ldots,x_\ell)=\prod_{k=1}^{\ell}(1-x_k+\alpha_\ell)
\]
for $i\in[m]$, $j\in[2^k]$, where $\alpha=\operatorname{bin}_\ell(j-1)$. By the construction $\deg(Q_j)=\ell$.
Let $P_{i,j}=Q_j(s_{i,1},s_{i,2},\ldots,s_{i,\ell})$.
Consider a tree-like $\mathrm{Th}(1)$-refutation of $\mathrm{PHP}_{2^\ell}^m$ of size $S$: $f_1\ge0,f_2\ge0,\ldots,f_S\ge0$, where $f_i$ are linear real polynomials over variables $p_{i,j}$ and $f_S\ge0$ is unsatisfiable on Boolean cube. For each of the inequalities on the following conditions hold:
\begin{enumerate}
\item[(a)] $f_i\ge0$ is semantically implied by $f_j\ge0$ and $f_k\ge0$ on the Boolean cube for $j,k<i$.
\item[(b)] $f_i$ is a linear representation of an axiom of $\mathrm{PHP}_{2^\ell}^m$;
\end{enumerate}
Let $F_i$ be a polynomial obtained of substitution $p_{j,k}:=P_{j,k}$ to $f_i$ for all $j\in[m]$; $k\in[2^\ell]$.
Consider a sequence of inequalities $F_1\ge0,\ldots,F_S\ge0$. Observe that $F_S\ge0$ is unsatisfiable on the Boolean cube since $P_{i,j}\in\{0,1\}$ on the Boolean cube. Let us verify that the sequence $F_1\ge0,\ldots,F_S\ge0$ may be extended to a correct tree-like $\mathrm{Th}(\ell)$ refutation of $\mathrm{BPHP}_{2^\ell}^m$:
\begin{enumerate}
\item[(a)] If $f_i\ge0$ is semantically implied by $f_j\ge0$ and $f_k\ge0$, then $F_i\ge0$ is also implied by $F_j\ge0$ and $F_k\ge0$, since $P_{i,j}$ is Boolean on the Boolean cube.
\item[(b)] If $f_i$ is a linear representation of a
\begin{itemize}
\item hole axiom then $f_i\ge0$ is equivalent to the function $(1-p_{a,b})+(1-p_{c,b})\ge1$ on $\{0,1\}^{\operatorname{Vars}(\mathrm{PHP}_{2^\ell}^m)}$ for $a,c\in[m]$, $b\in[2^\ell]$. Thus $F_i\ge0$ is also equivalent to $(1-P_{a,b})+(1-P_{c,b})\ge1$ on the Boolean cube. Observe that the restriction of $(1-P_{a,b})+(1-P_{c,b})\ge1$ to the Boolean cube coincides with the predicate $s_a\ne\operatorname{bin}_\ell(b)\vee s_c\ne\operatorname{bin}_\ell(b)$ which is an axiom of $\mathrm{BPHP}_{2^\ell}^m$.
\item pigeon axiom then $f_i\ge0$ is equivalent to $\sum_{j=1}^{2^\ell}p_{a,j}\ge1$ on the Boolean cube for some $a\in[m]$. Thus $F_i\ge0$ is equivalent to $\sum_{j=1}^{2^\ell}P_{a,j}\ge1$ on $\{0,1\}^{\operatorname{Vars}(\mathrm{BPHP}_{2^\ell}^m)}$. Observe that the latter inequality is identically true, since $P_{a,j}$ is equivalent to $s_a=\operatorname{bin}_\ell(j-1)$, so for exactly one value of $j\in[2^\ell]$, $P_{a,j}=1$. Since $F_i\ge0$ is identically true it can be semantically derived from two arbitrary axioms of $\mathrm{BPHP}_{2^\ell}^m$.
\end{itemize}
\end{enumerate}
It is easy to see that the size of the resulting refutation is at most $3S$.
\end{proof}
\begin{authorclaim}{Proposition44 ([4])}
For $m>n$ there is a tree-like Cutting Planes (a subsystem of $\mathrm{Th}(1)$) refutation of $\mathrm{PHP}_n^m$ of size $O(m^2n)$.
\end{authorclaim}
\begin{authorclaim}{Proposition6}
For $m>2^\ell$ there is a tree-like $\mathrm{Th}(\ell)$ refutation of $\mathrm{BPHP}_{2^\ell}^m$ of size $O(m^2 2^\ell)$.
\end{authorclaim}
\begin{proof}
Follows from Propositions43 and44.
\end{proof}
